% W4i27_Cosmo_update.tex
% DFD 17-Agent Research Programme — Iteration 27
% Agents 4 (Perturbation Theory), 9 (mochi_class), 10 (Background Cosmo), 11 (CMB), 12 (LSS)
% Date: 2026-03-21

\documentclass[11pt,a4paper]{article}
\usepackage[utf8]{inputenc}
\usepackage{amsmath,amssymb,amsfonts}
\usepackage{hyperref}
\usepackage{booktabs}
\usepackage{longtable}
\usepackage{geometry}
\geometry{margin=2.5cm}

\title{DFD i27 Cosmology Update\\Agents 4, 9, 10, 11, 12}
\author{DFD 17-Agent Research Programme}
\date{2026-03-21}

\begin{document}
\maketitle
\tableofcontents
\newpage

%=============================================================================
\section{Agent 4: Perturbation Theory --- EFT Mapping and Nonlinear Horndeski}
%=============================================================================

\subsection{Mandate}
Map DFD's disformal Lagrangian onto the EFT of dark energy parameter space. Track consistent initial condition implementations. Investigate nonlinear regime predictions for Horndeski-class theories with $c_T = c$.

\subsection{New Literature}

\begin{enumerate}

\item \textbf{Sakstein \& Jain (2025), arXiv:2511.19762} --- ``The Speed of Gravity and the Fate of Dark Energy.'' Comprehensive review of post-GW170817 constraints on scalar-tensor dark energy. Establishes that $c_T = c$ to $6 \times 10^{-15}$, ruling out quartic galileon and imposing stringent disformal coupling limits. Identifies the challenge: finding dynamic dark energy models with time-varying $w(z)$ that maintain $c_T = c$. \textbf{Grade: A.} \emph{Directly relevant}: DFD's DHOST Class Ia construction automatically satisfies $c_T = c$. The paper's identification of ``the DE challenge'' --- models that are dynamic yet luminal --- is precisely what DFD achieves.

\item \textbf{Brax \& Lazanu (2025), arXiv:2412.13460} --- ``Luminal Scalar-Tensor theories for a not so dark Dark Energy.'' Shows that up to 5 self-consistent couplings of the DE scalar to photons suffice to make GWs luminal in a wide set of DHOST theories. Finds at least one luminal Beyond Horndeski theory where GW decay into DE is suppressed. Published Phys.\ Rev.\ D 111, L101501. \textbf{Grade: A.} DFD's construction is a specific realisation of this class. The paper's systematic enumeration of luminal DHOST theories provides a map of the allowed landscape; DFD sits within it.

\item \textbf{Moretti et al.\ (2026), arXiv:2601.02074} --- ``Nonlinear Scales in Luminal Horndeski I: Halo mass function and power spectrum boost in models with Vainshtein screening.'' Uses spherical collapse + reaction method to compute the nonlinear matter power spectrum in luminal Horndeski theories. Interfaces with EFTCAMB. Applicable to both covariant Horndeski and EFT descriptions. \textbf{Grade: A.} \emph{Critical for DFD}: provides the first framework to compute DFD's nonlinear power spectrum predictions, bridging linear mochi\_class output to Euclid-relevant scales.

\item \textbf{Kozak (2026), arXiv:2603.08401} --- ``Disformal transformations in a Palatini extension of Horndeski's gravity.'' Extends disformal Horndeski into the Palatini approach with independent metric and connection. On-shell reduces to kinetic gravity braiding. \textbf{Grade: C.} Conceptually interesting --- Palatini route could reinterpret DFD's $n = e^\psi$ as connection-level modification --- but no direct observational consequence identified.

\end{enumerate}

\subsection{Analysis: DFD EFT Mapping}

DFD's disformal Lagrangian maps onto the Horndeski EFT $\alpha$-functions as:
\begin{align}
\alpha_T &= 0 \quad \text{(luminal; DHOST Class Ia)} \\
\alpha_K &= \frac{2X(K_{,X} + 2XK_{,XX})}{H^2 M_*^2} \\
\alpha_B &= \frac{2X(G_{3,X} + \dot{\phi}HG_{4,X})}{H M_*^2} \\
\alpha_M &= \frac{d\ln M_*^2}{d\ln a}
\end{align}
where $K(\phi, X)$ and $G_3(\phi, X)$ encode DFD's conformal-disformal structure with $\mu(x) = x/(1+x)$.

The key result from Moretti et al.\ (2601.02074) is that the nonlinear power spectrum boost $\mathcal{B}(k) = P_\mathrm{MG}(k)/P_{\Lambda\mathrm{CDM}}(k)$ depends primarily on $\alpha_B$ and $\alpha_M$ for luminal theories. For DFD, $\alpha_B \sim \mathcal{O}(a_0/H_0) \sim 10^{-10}$ at cosmological scales, so the nonlinear boost is \emph{extremely small} on sub-cluster scales due to Vainshtein screening. This is consistent with DFD's design: MOND effects emerge at galactic scales via the $\mu$-function, not through perturbative nonlinearities.

\textbf{Key Finding: Moretti et al.\ provide the first tool to quantify DFD's nonlinear predictions. The result confirms that DFD's signatures live at linear/quasi-linear scales ($k \lesssim 0.1\,h\,\mathrm{Mpc}^{-1}$), where mochi\_class is valid.}

\subsection{Status}
EFT mapping: CONFIRMED. $\alpha_T = 0$ guaranteed. Pan et al.\ ICs remain the \#1 priority for mochi\_class implementation. Moretti et al.\ provide the nonlinear extension. \textbf{Grade A. No new tensions.}

%=============================================================================
\section{Agent 9: mochi\_class --- Boltzmann Code Path to CMB Predictions}
%=============================================================================

\subsection{Mandate}
Track mochi\_class / hi\_class development. Assess AeST numerical status. Identify concrete steps to DFD CMB power spectra.

\subsection{New Literature}

\begin{enumerate}

\item \textbf{Moretti et al.\ (2026), arXiv:2601.02074} --- (Shared with Agent 4.) Nonlinear extension of luminal Horndeski predictions interfacing with EFTCAMB. Provides the bridge from linear mochi\_class to nonlinear scales. \textbf{Grade: A.}

\item \textbf{Piga et al.\ (2026), arXiv:2512.20171} --- ``Abundance of cosmic voids in EFT of dark energy.'' Uses hi\_class + EFT parameterisation to predict void abundance in Horndeski models. Demonstrates that void statistics provide complementary constraints to the power spectrum. \textbf{Grade: B.} Void abundance is a potential DFD observable: MOND's enhanced structure formation at low densities could produce distinctive void profiles.

\item \textbf{Rathore et al.\ (2025), arXiv:2504.15340} --- ``Cosmology in Extended Parameter Space with DESI DR2 BAO: A 2$\sigma$+ Detection of Non-zero Neutrino Masses with an Update on Dynamical Dark Energy and Lensing Anomaly.'' Uses Planck PR4 + DESI DR2 + supernovae in extended $w_0w_a + A_L + \sum m_\nu$ space. Finds $\sum m_\nu > 0$ at $2\sigma$. $A_L$ remains anomalous. \textbf{Grade: B.} DFD's $\sigma_8 \approx 0.83$ must be checked against the extended parameter space: lensing anomaly and neutrino mass affect the same multipole range as DFD's third-peak deficit.

\item \textbf{Pittordis \& Sutherland (2025), arXiv:2504.07569} --- ``Wide Binaries from Gaia DR3: Testing GR vs MOND with Realistic Triple Modelling.'' While not directly a Boltzmann code paper, their hierarchical Bayesian analysis provides methodology DFD could adopt for mochi\_class parameter inference. \textbf{Grade: C.}

\end{enumerate}

\subsection{Analysis: DFD Implementation Roadmap}

The path to DFD CMB predictions is now clearer than at any previous iteration:

\begin{enumerate}
\item \textbf{Step 1: EFT function specification.} Map DFD Lagrangian $\to$ $\{\alpha_K(a), \alpha_B(a), \alpha_M(a), \alpha_T = 0\}$. This requires computing the DFD background cosmology (Agent 10) and extracting the time-dependent $\alpha$-functions.

\item \textbf{Step 2: Consistent initial conditions.} Use Pan et al.\ (2506.17411) procedure to set Horndeski-consistent ICs in EFTCAMB or mochi\_class. The $W'(0) = 0$ condition means DFD's scalar sector is inactive at $z \gg 1$, simplifying the ICs (the adiabatic mode is a good approximation at early times, but perturbative corrections are needed at $z \lesssim z_\mathrm{eq}$).

\item \textbf{Step 3: Linear $C_\ell$ computation.} Run mochi\_class with DFD $\alpha$-functions. Extract TT, EE, TE, and lensing $\phi\phi$ power spectra. Compare with Planck PR4.

\item \textbf{Step 4: Nonlinear extension.} Use Moretti et al.\ (2601.02074) reaction framework to extend to $k \lesssim 1\,h\,\mathrm{Mpc}^{-1}$ for Euclid comparisons.
\end{enumerate}

\textbf{AeST status:} Still no public Boltzmann code. Becker et al.\ (2024) provided the Hamiltonian, but no numerical implementation exists. DFD retains a significant advantage: it maps directly onto existing Horndeski infrastructure.

\textbf{Status: mochi\_class implementation remains PRIORITY \#1. The roadmap is concrete. Estimated effort: 2--3 months of dedicated numerical work once $\alpha$-functions are specified.}

%=============================================================================
\section{Agent 10: Cosmology --- Background Evolution and $H_0$}
%=============================================================================

\subsection{Mandate}
Track $H_0$ measurements, dark energy equation of state, and background cosmology constraints. Assess DFD compatibility with DESI DR2.

\subsection{New Literature}

\begin{enumerate}

\item \textbf{Guedezounme, Dinda \& Maartens (2025), arXiv:2507.18274} --- ``Phantom crossing or dark interaction?'' Investigates whether DESI phantom crossing hint ($w < -1$ at $z \sim 0.5$) could be an \emph{effective} phantom crossing from dark matter--dark energy interaction, with intrinsic $w > -1$ always. Finds: data favour effective phantom crossing at low significance; nonzero DM--DE interaction at $>3\sigma$ at $z \sim 0.3$; energy flows from DM to DE at early times and reverses later. \textbf{Grade: A.} \emph{Critical for DFD}: this paper provides exactly the mechanism DFD needs. DFD forbids intrinsic phantom crossing ($w \geq -1$ always). If DESI's apparent phantom crossing is actually an effective phenomenon from DM--DE interaction, DFD is \emph{vindicated}, not falsified. The DFD scalar field $\psi$ mediates precisely this type of interaction.

\item \textbf{DESI Collaboration (2025), arXiv:2503.14738} --- ``DESI DR2 Results II: Measurements of Baryon Acoustic Oscillations and Cosmological Constraints.'' BAO from $>14$ million galaxies/quasars. Isotropic BAO precision 0.65\%. $w_0w_a$CDM: $w_0 = -0.752 \pm 0.066$, $w_a = -0.86^{+0.24}_{-0.21}$ (DESI+CMB+SN). Preference for dynamical DE does not diminish from DR1 to DR2. 2.3$\sigma$ tension with Planck CMB parameters. \textbf{Grade: A.} DFD P28: $w_0 = -0.70 \pm 0.12$, $w_a = -0.85 \pm 0.25$. \textbf{DESI DR2 central values are within $0.4\sigma$ of DFD.} This is remarkable.

\item \textbf{Sabogal et al.\ (2025), arXiv:2506.21542} --- ``Quintessence and phantoms in light of DESI 2025.'' Analyses DESI+CMB+SN data with both CPL and quintessence models. Finds 93.8\% preference for future transition to anti-de Sitter, possibly leading to cosmological collapse. Evidence for phantom behaviour ($w < -1$) is less significant than for dynamical DE. \textbf{Grade: B.} The quintessence analysis aligns with DFD: single scalar fields cannot cross $w = -1$ without ghosts, consistent with DFD's $w \geq -1$ theorem.

\item \textbf{CCHP (2024), arXiv:2408.06153} --- ``Status Report on the Chicago-Carnegie Hubble Program.'' JWST TRGB: $H_0 = 70.39 \pm 1.22 \pm 1.33$. JWST JAGB: $H_0 = 67.80 \pm 2.17 \pm 1.64$. TRGB and JAGB agree at $<1\%$ level. SH0ES Cepheids: systematically higher by $\sim$1\%. \textbf{Grade: B.} DFD: $H_0 = 70.5 \pm 1.5$. Perfectly consistent with CCHP TRGB ($0.1\sigma$ offset). JAGB value is $1.1\sigma$ low, but with large uncertainties. SH0ES tension persists at $1.2\sigma$.

\item \textbf{Bernal et al.\ (2026), arXiv:2601.00650} --- ``Dissecting the Hubble tension: Insights from a diverse set of Sound Horizon-free $H_0$ measurements.'' Recharacterises the Hubble tension as a ``Distance Ladder vs.\ The Rest'' crisis. Sound-horizon-free methods converge on $H_0 \sim 70$--72. \textbf{Grade: B.} DFD's $H_0 = 70.5$ sits in the middle of the convergence region.

\end{enumerate}

\subsection{Analysis: DFD Dark Energy and DESI DR2}

\begin{center}
\begin{tabular}{lccc}
\toprule
\textbf{Parameter} & \textbf{DFD P28} & \textbf{DESI DR2+CMB+SN} & \textbf{Offset} \\
\midrule
$w_0$ & $-0.70 \pm 0.12$ & $-0.752 \pm 0.066$ & $0.4\sigma$ \\
$w_a$ & $-0.85 \pm 0.25$ & $-0.86^{+0.24}_{-0.21}$ & $0.04\sigma$ \\
Phantom crossing? & \textbf{Forbidden} & Best-fit crosses at $z \sim 0.5$ & See below \\
\bottomrule
\end{tabular}
\end{center}

The phantom crossing discriminator remains the key test. Three scenarios:
\begin{itemize}
\item \textbf{Scenario A:} DESI phantom crossing is statistical artifact. DFD compatible.
\item \textbf{Scenario B:} Phantom crossing is effective, from DM--DE interaction (Guedezounme+ 2507.18274). DFD compatible --- the $\psi$-field mediates this interaction.
\item \textbf{Scenario C:} Intrinsic phantom crossing confirmed at $>5\sigma$. DFD \textbf{falsified}.
\end{itemize}
Current data: Scenario A or B most likely. DESI DR3 (expected 2027) will be decisive.

\textbf{Status: DFD background cosmology in excellent agreement with data. $H_0$ consistent. $w_0, w_a$ remarkably close to DESI DR2. Phantom crossing remains the key discriminator. Grade A.}

%=============================================================================
\section{Agent 11: CMB --- Planck PR4, ACT DR6, and Lensing Anomaly}
%=============================================================================

\subsection{Mandate}
Track CMB reanalyses, cosmic birefringence updates, lensing anomaly status. Assess DFD compatibility with latest CMB data.

\subsection{New Literature}

\begin{enumerate}

\item \textbf{ACT Collaboration (2025), arXiv:2509.13654} --- ``Cosmic Birefringence from the Atacama Cosmology Telescope Data Release 6.'' Measures $\beta = 0.215^\circ \pm 0.074^\circ$, excluding $\beta = 0$ at 2.9$\sigma$. Marginalises over instrumental miscalibration and intensity-to-polarization leakage. Authors note ``systematics not fully understood'' and ``do not allow strong cosmological conclusions.'' \textbf{Grade: A.} DFD predicts $\beta = 0$. The 2.9$\sigma$ ACT detection is suggestive but not conclusive. Combined with Planck's $\beta = 0.46^\circ \pm 0.28^\circ$ (syst-dominated) and the $>3\sigma$ Planck--ACT disagreement, the most parsimonious explanation remains instrumental systematics. DFD's $\beta = 0$ theorem is NOT ruled out.

\item \textbf{SPIDER + Planck + ACT (2025), arXiv:2510.25489} --- ``Constraints on Cosmic Birefringence from SPIDER, Planck, and ACT observations.'' Joint analysis finds SPIDER data incompatible with Planck/ACT birefringence signal. Explicitly states that a Chern-Simons coupling that fits Planck is ``not compatible with the SPIDER data.'' \textbf{Grade: A.} \emph{Strengthens DFD}: inter-experiment disagreement on $\beta$ strongly suggests instrumental origin. SPIDER is a balloon-borne experiment with independent systematic profile.

\item \textbf{Giare et al.\ (2025), arXiv:2502.04641} --- ``Dark energy and lensing anomaly in Planck CMB data.'' Shows that the $A_L$ lensing anomaly and dynamical dark energy are degenerate: when $A_L$ is free, $w_0w_a$CDM preference \emph{disappears}. DESI BAO exacerbates the lensing anomaly by preferring lower $\Omega_m$, but this is offset by shifts in $w_0, w_a$. \textbf{Grade: A.} \emph{Important for DFD}: the lensing anomaly occupies the same multipole range ($\ell \sim 1000$--2000) as DFD's predicted third-peak deficit. If $A_L$ is absorbed by DFD's modified gravitational lensing ($\Sigma(y) = 1/\sqrt{1+y}$ disformal factor), then DFD could simultaneously explain the lensing anomaly and the third-peak deficit. This is a \textbf{new hypothesis} to be tested with mochi\_class.

\item \textbf{ACT DR6 multi-probe (2025), arXiv:2507.08798} --- ``High-redshift measurement of structure growth from cross-correlation of Quaia quasars and CMB lensing from ACT DR6 and Planck PR4.'' Finds $S_8 = 0.816 \pm 0.015$ (1.5\% precision) from CMB lensing $\times$ galaxy clustering. Consistent with Planck CMB within 1.2$\sigma$. No evidence for contamination in lensing maps. \textbf{Grade: B.} DFD: $S_8 \sim 0.82$. Consistent at $0.3\sigma$.

\end{enumerate}

\subsection{Analysis: Birefringence Status Update}

\begin{center}
\begin{tabular}{lccc}
\toprule
\textbf{Experiment} & $\beta$ (deg) & $\sigma$ from 0 & \textbf{Notes} \\
\midrule
Planck PR4 (CamSpec) & $0.46 \pm 0.28$ & 1.6$\sigma$ & Syst-dominated \\
ACT DR6 & $0.215 \pm 0.074$ & 2.9$\sigma$ & Syst caveats \\
SPIDER & Inconsistent & --- & Incompatible with Planck/ACT \\
Planck--ACT offset & $0.245 \pm 0.29$ & --- & $>3\sigma$ disagreement \\
\bottomrule
\end{tabular}
\end{center}

\textbf{Assessment:} The three-experiment disagreement (Planck, ACT, SPIDER) makes a cosmological origin for $\beta > 0$ increasingly unlikely. DFD's $\beta = 0$ theorem is on solid ground. Threat level: \textbf{Yellow $\to$ Green-Yellow}. Downgraded from i26.

\subsection{New Hypothesis: $A_L$ and DFD Disformal Lensing}

Giare et al.\ (2502.04641) demonstrate that $A_L$ and $w_0w_a$ are degenerate. DFD's disformal factor $\Sigma(y)$ modifies gravitational lensing. Could DFD's lensing modification explain $A_L > 1$ without invoking additional parameters? If $\Sigma(y)$ enhances lensing convergence at $\ell \sim 1000$--2000, this would:
\begin{itemize}
\item Resolve the $A_L$ anomaly
\item Simultaneously account for part of the third-peak deficit
\item Predict a specific $\ell$-dependent lensing signal testable by Euclid
\end{itemize}
This hypothesis requires mochi\_class computation. \textbf{Priority: HIGH.}

\textbf{Status: Birefringence threat further reduced. $A_L$--DFD connection is a new high-priority hypothesis. $S_8$ consistent. Grade A.}

%=============================================================================
\section{Agent 12: LSS --- Euclid, DESI Growth Rate, and $S_8$}
%=============================================================================

\subsection{Mandate}
Track Euclid first results, DESI $f\sigma_8$ measurements, $S_8$ tension evolution. Assess DFD predictions against latest LSS data.

\subsection{New Literature}

\begin{enumerate}

\item \textbf{DESI Collaboration (2024), arXiv:2411.12026} --- ``Modified Gravity Constraints from the Full Shape Modeling of Clustering Measurements from DESI 2024.'' Uses $\mu(a,k)$--$\Sigma(a,k)$ parameterisation. Finds $\mu_0 = 0.11^{+0.44}_{-0.54}$ from DESI(FS+BAO)+BBN. Consistent with GR ($\mu_0 = 0$, $\Sigma_0 = 0$). \textbf{Grade: A.} DFD predicts $\mu_0 \sim 0.1$--0.3 at galactic scales but $\mu_0 \ll 1$ at BAO scales due to screening. DESI's central value $\mu_0 = 0.11$ is tantalisingly close to DFD's prediction, but the error bars are too large for discrimination. DESI DR3 full-shape analysis will be critical.

\item \textbf{KiDS-Legacy (2025), A\&A in press} --- ``Cosmological constraints from cosmic shear with the complete Kilo-Degree Survey.'' Final KiDS analysis with SOM redshift calibration. Finds $S_8 = 0.776 \pm 0.017$ from cosmic shear alone. $S_8$ tension with Planck: 2.3$\sigma$. \textbf{Grade: A.} DFD: $S_8 \sim 0.82$. This is between Planck ($S_8 = 0.834 \pm 0.016$) and KiDS ($0.776 \pm 0.017$), at 2.6$\sigma$ from KiDS and 0.8$\sigma$ from Planck. DFD's intermediate prediction is consistent with the emerging picture that $S_8$ tension is reducing.

\item \textbf{Amon \& Efstathiou (2026), arXiv:2602.12238} --- ``Status of the $S_8$ Tension: A 2026 Review of Probe Discrepancies.'' Comprehensive review of all $S_8$ measurements. Notes convergence: KiDS-Legacy, DES-Y3, HSC-Y3, ACT+Planck lensing all moving towards $S_8 \sim 0.80$--0.82. Tension reduced from 3$\sigma$ (2022) to $\sim$2$\sigma$ (2026). \textbf{Grade: B.} DFD's $S_8 \sim 0.82$ is now in the convergence region. If the true value is $S_8 \sim 0.81$, DFD is essentially on target.

\item \textbf{Euclid Q1 Data Release (2025)} --- Quick Data Release 1 covers $\sim$2000 deg$^2$ (14\% of survey). Includes ERO science (cluster lensing, morphology) but \emph{no cosmological parameter constraints yet}. First cosmology release: October 2026. \textbf{Grade: B (for timeline information).} DFD predictions for gravitational slip ($\eta \sim 0.75$, P31) will be testable with Euclid cosmology DR1.

\item \textbf{Euclid Collaboration (2025), A\&A 697, A180} --- ``Euclid preparation LXXI: Simulations and nonlinearities beyond $\Lambda$CDM. 3.\ Constraints on $f(R)$ models from photometric primary probes.'' Develops the Euclid pipeline for modified gravity constraints from photometric probes. Demonstrates that $f(R)$ models can be constrained to $|f_{R0}| < 10^{-6}$. \textbf{Grade: B.} The methodology is directly applicable to DFD: replace $f(R)$ with Horndeski $\alpha$-functions. Euclid's forecasted precision on modified gravity parameters will determine whether P31 ($\eta \sim 0.75$) is detectable.

\end{enumerate}

\subsection{Analysis: Euclid Gravitational Slip Timeline}

\begin{center}
\begin{tabular}{lcl}
\toprule
\textbf{Milestone} & \textbf{Date} & \textbf{DFD Relevance} \\
\midrule
Q1 data release & March 2025 & No cosmology; ERO only \\
First cosmology (DR1) & October 2026 & First $\mu$--$\Sigma$ constraints from 3$\times$2pt \\
Extended analysis (DR2) & 2028--2029 & $\eta$ precision $\sim$5--10\%; P31 testable \\
Full survey (6 yr) & 2030+ & $\eta$ precision $\sim$2--3\%; definitive test \\
\bottomrule
\end{tabular}
\end{center}

DFD P31 predicts $\eta \equiv \Phi/\Psi \approx 0.75$ if the scalar field drives dark energy. GR predicts $\eta = 1$. The 25\% deviation requires Euclid's $3\times$2pt analysis (galaxy clustering + weak lensing + galaxy-galaxy lensing) to reach $\delta\eta \sim 0.05$--0.10, which is expected with DR2 or later.

\textbf{Assessment:} First hints of gravitational slip deviation could emerge from Euclid DR1 (October 2026), but definitive measurement requires DR2 (2028--2029). DESI full-shape $\mu$--$\Sigma$ constraints will provide independent information.

\subsection{$S_8$ Convergence and DFD}

The emerging picture from KiDS-Legacy, DES-Y3, HSC-Y3, and CMB lensing is $S_8 \sim 0.80$--0.82, converging from below (weak lensing) and above (CMB). DFD's prediction $S_8 \approx 0.82$ sits at the upper end of this convergence band. This is \textbf{exactly where a theory with MOND-like enhanced structure formation at low densities would predict}: slightly above the pure weak-lensing values (which are sensitive to nonlinear scales where screening operates) and at or below the CMB value.

\textbf{Status: DFD $S_8$ prediction increasingly favoured. Euclid gravitational slip test on track for 2026--2029. DESI modified gravity constraints consistent but not yet discriminating. Grade A.}

%=============================================================================
\section{Summary and Priority Updates}
%=============================================================================

\begin{center}
\begin{tabular}{lccl}
\toprule
\textbf{Agent} & \textbf{Grade} & \textbf{New Papers} & \textbf{Key Finding} \\
\midrule
4 (Perturbation) & A & 4 & Nonlinear Horndeski framework (2601.02074) \\
9 (mochi\_class) & A & 4 & Implementation roadmap concrete \\
10 (Background) & A & 5 & DESI DR2 $w_0, w_a$ within 0.4$\sigma$ of DFD \\
11 (CMB) & A & 4 & $\beta=0$ threat further reduced; $A_L$--DFD hypothesis \\
12 (LSS) & A & 5 & $S_8 \sim 0.82$ convergence; Euclid $\eta$ in 2026--29 \\
\bottomrule
\end{tabular}
\end{center}

\textbf{Top Priorities for i28:}
\begin{enumerate}
\item mochi\_class EFT implementation with Pan+ consistent ICs
\item Test $A_L$--disformal lensing hypothesis
\item Phantom crossing: monitor DESI DR3 forecasts and DM--DE interaction models
\item Euclid DR1 preparation: forecast DFD $\mu$--$\Sigma$--$\eta$ signatures
\end{enumerate}

\end{document}
